%----------------------------------------------------------------------------
\chapter{A feladat értelmezése}
%----------------------------------------------------------------------------

\section{ A feladat célja és értelmezése}
A feladat fő célja egy autonóm ágensalapú rendszer megvalósítása, amely képes erőforrás-allokációk
értékelésére, javaslatok megfogalmazására önállóan, majd felhasználói jóváhagyással azok végrehajtására.

A beavatkozási javaslat készítése egy külső nagy nyelvi modell segítségével történik,
a futó alkalmazásokból kinyert metrikai adatok és korlátok elküldésével. A jóváhagyott változtatások végrehajtását a
rendszer szintén egy nagy nyelvi modellel és annak eszközhasználatával végzi. A tényleges módosításokat az alkalmazás 
által biztosított, biztonsági korlátokkal ellátott eszközök végzik.

A működés vizsgálatához különböző terhelési és erőforrás-beállításokkal futó mintaalkalmazásokat definiáltam, amelyek
segítségével túlméretezett, processzorkorlátozott és memóriakockázatos helyzeteket,
valamint megfelelően méretezett kontrollesetet is vizsgálok.

A feladatkiírásban több lehetséges vizsgált adatforrás és beavatkozási módszer szerepel.
Az én applikációm ezek közül a processzor terhelési és memóriahasználathoz kapcsolódó adatok alapján
az optimalizálandó alkalmazások erőforrás-beállításainak és példányszámainak megfelelő módosítására
törekszik. Ez lehetővé teszi mind a vertikális, mind az alkalmazásszintű horizontális skálázási döntések
vizsgálatát egy jól elkülönített próbakörnyezetben.

Itt fontos megjegyezni, hogy a felhőszolgáltatónál üzembe helyezett rendszer optimalizálási változtatásai
a bérelt virtuális gépek számát vagy méretét nem érintik, ezért
a változtatások se nem csökkentik, se nem növelik a bérelt gépek költségét.

\section{A vizsgált metrikák és a beavatkozások hatóköre}
A CPU (Central Processing Unit) -használat azt vizsgálja, hogy az applikáció mekkora processzorkapacitást
vesz igénybe. Itt a vizsgált mennyiség mértékegysége a millicore, amely egy CPU-mag
kapacitásának az egy ezredét jelenti. A processzorhasználatot adott időintervallumra kiátlagoljuk, ez adja az átlagos
CPU-használatot, illetve a fontosabb metrika a használat 95. percentilise, azaz p95 értéke. Ez az érték az, amelyet a
vizsgált minták 95 százaléka nem lép át. Az átlaghoz képest jobban mutatja a magasabb terhelésű időszakokat, ezért ezt használjuk
a javaslatokhoz.\\

A CPU korlátozás egy olyan jelenség, ami akkor lép fel,
ha az adott applikáció több CPU-t szeretne használni, mint amennyi a rá kiszabott maximum limit,
ezért az alkalmazás működését a rendszer visszatartja, ami lassíthatja a működést.
A mi esetünkben a jelenség fellépésének százalékos arányát mérjük a vizsgált időszakon belül.
Azt, hogy a periódusok hány százalékában történt korlátozás.\\

A memóriahasználat az applikáció által használt memória mennyiségét jellemzi.
A rendszer itt a memóriacsúcsot veszi figyelembe, amely a vizsgált időszakban
fellépő legnagyobb memóriahasználat értéke. Amennyiben ez megközelíti a kiszabott
memórialimit értékétét, az kockázatot jelent.\\

A beavatkozások hatókörébe tartoznak mind a vertikális, mind a horizontális skálázási döntések,
ezek vertikális szempontból kimerülnek az alkalmazáspéldányok CPU- és memóriaigényeinek,
illetve ezek korlátainak csökkentésével vagy növelésével, horizontális szempontból pedig
az alkalmazások példányszámainak módosításával.\\

A rendszer a megfigyelést és a javaslatok generálását automatikusan végzi. A javaslatokat és a hozzájuk tartozó
mérési adatokat, illetve a végrehajtások eredményeit eltárolja és megjeleníthetővé teszi.
A végrehajtáshoz a felhasználónak manuálisan kell jóváhagynia a javasolt változtatásokat.

\section{Funkcionális követelmények}
\begin{enumerate}
    \item A rendszer gyűjtse és jelenítse meg a megfigyelt alkalmazások állapotát és metrikáit.
    \item A mérési adatok és aktuális erőforrás-beállítások alapján ismerje fel a javaslatot igénylő helyzeteket.
    \item Adjon pontos erőforrás vagy példányszám változtatási javaslatokat és jelenítse meg őket.
    \item Adjon lehetőséget a javaslatok elfogadására és elutasítására.
    \item Az elfogadott javaslatokban szereplő változtatásokat hajtsa végre.
    \item A sikeresen végrehajtott válzotzatásokhoz biztosítson lehetőséget a változtatások visszaállítására.
    \item A javaslatokat, végrehajtásokat és visszavonásokat tárolja, adjon lehetőséget azok megjelenítésére.
\end{enumerate}

\section{Kiértékelés módja}
A kiértékelésnél 5 forgatókönyvet alapján vizsgálunk, amelyeket 4 mintalkalmazáson fogunk kiértékelni.
A forgatókönyvek a következők:
\begin{itemize}
    \item \emph{cpu-hungry - vertikális skálázás} - Azt az esetet mutatja be, amikor az applikáció több CPU-t használna, mint ami allokálva van neki.
    Itt jellemzően a korlátozás közel 100 százalékos értéket érhet el, itt a CPU allokáció növelését várjuk.
    \item \emph{cpu-hungry - horizontális skálázás} - Ennél ugyan úgy a processzor korlátozott alkalmazást vizsgáljuk,
    itt párhuzamos kérésekkel terheljük az appot, ezért a példányszám növelését várjuk javaslatként.
    \item \emph{idle-app} - Ebben az esetben az allokált CPU jóval meghaladja azt a mennyiséget amelyre ennek az appnak szüksége van,
    ezért itt jellemzően a CPU kihasználtság nagyon alacsony. Itt a CPU allokáció csökkentését várjuk javaslatként.
    \item \emph{mem-cache} - Itt a memóriacsúcs megközelíti a kiszabott memória limitet, ezért fellép a memória 
    elfogyásának kockázata. A várható javaslat a memória limit növelése. 
    \item \emph{web-api} - Ez egy normál kontrollként használt eset. A CPU kihasználtsága nem túl magas, de nem is alacsony,
    a memóriából is jelentős rész használva van, de nem közelíti meg a memóriacsúcs a limitet.
    Itt arra számítunk, hogy nem kapunk változtatási javaslatot.
\end{itemize}

A kiértékelésnél a helyességet úgy fogjuk vizsgálni, hogy megnézzük, tényleg megfelelő javaslatot ad a rendszer,
illetve hogy végrehajtásnál tényleg a javaslatban szereplő erőforrás-allokációk, illetve példányszám változások lépnek érvénybe.




